\begin{textsample}{POS Dim 2 – sp – Score 20.00 – sp/sp\_inf\_finalidade.txt}  \label{ex:f2_pos_204}
Identidade e finalidade da educação infantil História da Educação Infantil no Estado de São Paulo Recuperar a história da Educação Infantil no Brasil contribui para compreender o que a Base Nacional Comum Curricular ( BNCC ) representa em termos de direito à criança para esta etapa da educação.
Saber de onde partimos, onde estamos e aonde queremos chegar possibilita traçar novos caminhos.
Nos anos 1930, já existiam no Brasil instituições públicas de proteção à criança, mas foi na década de 1940 que as ações governamentais se efetivaram com um foco na filantropia, no higienismo, na puericultura, vertentes vinculadas às questões de cuidado e saúde às crianças.
As creches eram planejadas com rotinas de triagem, lactário, enfermagem e preocupação com a higiene do ambiente físico ( OLIVEIRA, 2005 ).
Com relação ao Estado de São Paulo, conforme Andrade ( 2010 ), as escolas maternais foram regularizadas em 1925 por meio de um decreto \textbf{estadual}.
Posteriormente, em 1935, quando Mário de Andrade dirigia o Departamento de Cultura do \textbf{município} de São Paulo, foram criados os parques infantis nos bairros operários que \textbf{atendiam} crianças de diferentes grupos etários em \textbf{horário} contrário ao da escola para atividades recreativas.
Nessa época, os interesses da burguesia, dos \textbf{trabalhadores} e do Estado fomentaram \textbf{políticas} públicas que regulamentaram o atendimento à infância.
Na década de 1950, a maioria das creches era de responsabilidade de sociedades filantrópicas e tinham como objetivo suprir as necessidades advindas da pobreza ( KUHLMANN, 2001 ).
Em 1964 foi criada a Fundação Nacional de Bem-Estar do Menor ( FUNABEM ), tendo como proposta uma educação compensatória, que buscava a redução do fracasso escolar das crianças das classes desfavorecidas.
O projeto continuou na década de 1970, com a atuação da Legião Brasileira de Assistência.
A primeira ação voltada à infância em âmbito \textbf{estadual} é promovida em 1966.
Sem abandonar totalmente os princípios higienistas e assistencialistas, é defendido no I Seminário sobre Creches no Estado de São Paulo o conceito de creche como “ um serviço que oferece um potencial capaz de garantir o desenvolvimento infantil, compensando as deficiências de um meio precário próprio das famílias de classe \textbf{trabalhadora} ” ( HADDAD, 1990, p.109 ).
Nesse evento, realizado pela Secretaria do Bem-Estar Social, a creche é apresentada como instituição de atenção à infância capaz de \textbf{atender} aos filhos da mãe \textbf{trabalhadora}, que tem como objetivo a promoção da família e a prevenção da marginalidade, mas quer sobretudo sensibilizar a sociedade civil para a \textbf{qualidade} do atendimento \textbf{ofertado} às crianças.
Buscando essa \textbf{qualificação}, a Secretaria passa a defender a necessidade de contar com \textbf{profissionais} especializados na área do desenvolvimento e educação infantil — do Serviço Social, da Psicologia, da Pedagogia e de outras áreas afins — para pensar e realizar o trabalho nas creches.
Contudo, influenciados pelo tecnicismo, esses \textbf{profissionais}, especialmente os do Serviço Social, mantêm um olhar \textbf{técnico} para o trabalho que prioriza as famílias mais do que as crianças.
Na década de 1970, com a promulgação da Lei nº 5.692, de 1971, uma das normativas \textbf{federais} define a função social da Educação Infantil e reconhece sua importância como etapa educacional, conforme se lê no \textbf{capítulo} 6, artigo 61, da referida lei : “ Os sistemas de ensino estimularão as empresas que tenham em seus serviços mães de menores de sete anos a organizar e manter, diretamente ou em cooperação, inclusive com o Poder Público, educação que preceda o ensino de 1º grau ”.
Em 1981, com a criação do Programa Nacional da Educação Pré-escolar, elaborado pelo MEC/COEPRE/Secretarias de Educação e pelo Mobral, observa -se um movimento inicial para a educação das infâncias, embora esta não estivesse ainda sendo tratada como força constitucional.
O Programa reconhecia a relevância de ações voltadas à infância frente ao impacto que esta tem no desenvolvimento do ser humano. > [... ] A educação pré-escolar é agora considerada como a primeira fase da educação, pois, estabelece a base de todo processo educativo, que consiste em a pessoa fazer -se progressiva e permanentemente \textbf{conquistando} -se a si mesma, integrando -se ao grupo social, delineando o seu presente e criando o seu futuro. ( BRASIL, 1981, p.5 ) Em São Paulo, a década de 1980 foi marcada por movimentos pró-creches que, influenciados pela luta das mulheres, apresentavam várias reivindicações aos poderes públicos.
Representando uma luta por direitos sociais e cidadania, tais movimentos resultaram na conquista da creche como um direito das crianças e da mulher \textbf{trabalhadora} ( MERISSE, 1997 ).
A Constituição Federal de 1988 ratifica à criança de 0 a 6 anos o direito de \textbf{frequentar} creches e pré-escolas.
Com a chegada da Lei de \textbf{Diretrizes} e Bases da Educação Nacional ( Lei nº 9.394 de 20 de dezembro de 1996 ), a Educação Infantil é integrada à Educação Básica.
Em 2006, a LDB passa por alterações e reduz o período da Educação Infantil para 0 a 5 anos em razão da ampliação do Ensino Fundamental para 9 anos.
Em 2013 é regulamentada a Lei nº 12.796/2013, que inclui na LDB a obrigatoriedade da matrícula de todas as crianças de 4 e 5 anos na Educação Infantil.
Nesse percurso histórico de avanços e conquistas da Educação Infantil brasileira, em dezembro de 2017 é \textbf{homologada} a Base Nacional Comum Curricular ( BNCC ), \textbf{atendendo} à Constituição Federal/88 e à LDB/96 e contemplando em seus princípios as \textbf{Diretrizes} Curriculares Nacionais da Educação Infantil ( DCNEI, Resolução CNE/CEB nº 5/2009 ).
Como desdobramento desse processo, a construção do \textbf{Currículo} Paulista para a Educação Infantil traz como premissas o binômio educar e cuidar, as interações e brincadeiras e a \textbf{garantia} dos direitos de aprendizagem e desenvolvimento das crianças — conviver, brincar, participar, explorar, expressar e conhecer -se, contempladas na BNCC.
No cenário \textbf{estadual}, de acordo com a Pesquisa Nacional por Amostra Domiciliar ( PNAD ), são \textbf{atendidas} aproximadamente 40 \% das crianças na creche e cerca de 93 \% das crianças na pré-escola, dados que apontam para a necessidade de \textbf{políticas} públicas voltadas a essa etapa da Educação Básica, como forma de atendimento à meta 01 do Plano Nacional de Educação ( PNE, de 25 de junho de 2014 ), que versa sobre a universalização da pré-escola e da ampliação na \textbf{oferta} de creche.
O Estado de São Paulo é constituído por uma população representativa de diversas regiões do país.
Tal especificidade evidencia a necessidade de se considerar a diversidade cultural no momento da construção do \textbf{Currículo} Paulista.
Como previsto na LDB, os \textbf{municípios} têm autonomia para definir as \textbf{políticas} públicas que viabilizem a \textbf{oferta} e o acesso a um atendimento de \textbf{qualidade}, de forma a respeitar o contexto social, histórico e cultural em que estão inseridos.
Nessa perspectiva, coube a esse \textbf{currículo} assegurar princípios para o atendimento à criança pequena nas creches e na pré-escola, instituições que devem acolhê -la e partilhar com sua família e/ou responsáveis os cuidados a que tem direito na infância — com seu corpo e pensamento, seus afetos e sua imaginação — e garantir as aprendizagens essenciais, respeitando a história construída no ambiente familiar e/ou na comunidade em que vive.

% matched lemmas: atender, capítulo, conquistar, currículo, diretriz, estadual, federal, frequentar, garantia, homologar, horário, município, oferta, ofertar, política, profissional, qualidade, qualificação, trabalhador, técnico
\end{textsample}
